% from ega3-7-fr.tex
\oldpage[III]{179}
Posons
\[
  T'(M)=T(A_s)\otimes_A M
\]
pour tout $A$-module à gauche $M$. Il existe une suite exacte
\[
  L'\longrightarrow L\longrightarrow M\longrightarrow0,
\]
où $L$ et $L'$ sont deux $A$-modules à gauche libres; comme $T$ et $T'$ sont
exacts à droite, on a donc le diagramme commutatif
\[
\xymatrix@C=22pt@R=25pt{
  T'(L')\ar[r]\ar[d]^{t_{L'}}
  &
  T'(L)\ar[r]\ar[d]^{t_L}
  &
  T'(M)\ar[r]\ar[d]^{t_M}
  &
  0
  \\
  T(L')\ar[r]
  &
  T(L)\ar[r]
  &
  T(M)\ar[r]
  &
  0
}
\]
où les deux lignes sont exactes; comme $t_L$ et $t_{L'}$ sont des
isomorphismes en vertu de (7.2.4), il en est de même de $t_M$ par le lemme des
cinq. Enfin, il est clair que $b)$ entraîne $b')$. Pour montrer que $b')$
entraîne $a)$, il suffit de prouver le lemme suivant.
\end{proof}

\begin{lemma}[7.2.5.1]
Soient $\mathcal K$, $\mathcal K'$ deux catégories abéliennes, $F$, $G$ deux
foncteurs covariants additifs de $\mathcal K$ dans $\mathcal K'$,
$f:F\to G$ un morphisme fonctoriel (T, I, 1.2) tel que, pour tout objet $E$
de la catégorie $\mathcal K$,
\[
  f_E:F(E)\longrightarrow G(E)
\]
soit un épimorphisme. Alors, si $F$ est exact à droite et $G$ semi-exact,
$G$ est exact à droite.
\end{lemma}

\begin{proof}
Tout revient en effet à démontrer que pour tout épimorphisme
$v:E'\to E$ dans $\mathcal K$,
\[
  G(v):G(E')\longrightarrow G(E)
\]
est un épimorphisme; or, on a le diagramme commutatif
\[
\xymatrix@C=35pt@R=30pt{
  F(E')\ar[r]^{F(v)}\ar[d]_{f_{E'}}
  &
  F(E)\ar[d]^{f_E}
  \\
  G(E')\ar[r]_{G(v)}
  &
  G(E)
}
\]
dans lequel $F(v)$, $f_{E'}$ et $f_E$ sont des épimorphismes; il en est donc
de même de $G(v)$.
\end{proof}

\begin{remark}[7.2.6]
Pour tout $A$-module à droite $N$, posons
\[
  T_N(M)=N\otimes_A M
\]
pour tout $A$-module à gauche $M$, de sorte que $T_N$ est un foncteur
covariant additif de $\mathbf{Ab}_A$ dans $\mathbf{Ab}$, exact à droite et
commutant aux sommes directes. Si on identifie canoniquement $T_N(A_s)$ à
$N$, on vérifie aussitôt que l'homomorphisme correspondant (7.2.2.1) devient
l'identité. On en conclut que le $A$-module à droite $N$ de l'énoncé de
(7.2.5, $c)$) est déterminé à un isomorphisme unique près et est
canoniquement isomorphe à $T(A_s)$. On peut encore dire que les morphismes
fonctoriels
\[
  T\longmapsto T(A_s),
  \qquad
  N\longmapsto T_N
\]
constituent une équivalence (T, I, 1.2) de la catégorie des $A$-modules à
droite et de la catégorie des foncteurs covariants additifs
$\mathbf{Ab}_A\to\mathbf{Ab}$ qui sont exacts à droite et commutent aux
sommes directes.
\end{remark}

\begin{proposition}[7.2.7]
Soient $A$ un anneau artinien à gauche, dont le quotient par son radical
$\mathfrak m$ est un corps $k$. Soit $T$ un foncteur covariant additif de
$\mathbf{Ab}_A$ dans $\mathbf{Ab}$, commutant aux sommes directes. Les
conditions de (7.2.5) sont alors aussi équivalentes à
\begin{enumerate}
  \item[d)] $T$ est semi-exact et l'homomorphisme
    \[
      T(\varepsilon):T(A_s)\longrightarrow T(k)
    \]
    déduit de l'homomorphisme canonique $\varepsilon:A_s\to k$ est surjectif.
\end{enumerate}
\end{proposition}

\begin{proof}
